\section{Заключение}
\label{sec:conclusion}

Спроектировано и реализовано вычисление LeaderRank на ориентированном невзвешенном графе, который не
помещается в память. Ключевые решения:
\begin{itemize}[itemsep=2pt]
  \item \textbf{Метрика --- LeaderRank.} Её ядро --- один потоковый проход по рёбрам
        ($A\T\cdot\mathrm{contrib}$), что и делает задачу разрешимой во внешней памяти.
  \item \textbf{Хранение --- CSC по приёмнику.} Указатели \texttt{sourcesPtr} ($\bigO(n)$) в ОЗУ,
        источники \texttt{sources} ($\bigO(m)$) на диске, чтение последовательное.
  \item \textbf{Параллелизм --- модель pull с разбиением по приёмнику.} Без блокировок и
        детерминированно при любом числе потоков.
  \item \textbf{Память --- ограничена числом вершин.} Лимит задаётся снаружи (\texttt{-Xmx}, cgroup),
        программа под него подстраивается, пик RSS измеряется и приводится.
\end{itemize}

Эти решения прямо отвечают критериям задачи: эффективность на больших графах, использование всех
ядер, устойчивость к гипер-узлам, воспроизводимость результата и чистота кода.

Главные числа. Параллельное ядро даёт ускорение $7{,}39\times$ на $24$ потоках, детерминированно
($38$ итераций при любом $P$). Реальные графы SNAP считаются out-of-core: web-Google ($5{,}1$ млн
рёбер) --- около $0{,}4$~с, soc-LiveJournal1 ($69$ млн рёбер) --- около $3{,}3$~с. Контрольный наивный
вариант под тем же лимитом памяти падает, а потоковый движок завершается --- это и подтверждает
ценность out-of-core подхода.
